\subsection{Revisão Histórica}

Na atual etapa da humanidade os sabonetes são produtos tão comuns que pode até mesmo passar despercebido que estamos trabalhando em sua formulação a mais de 5 mil anos. O primeiro registro de materiais parecidos com o sabão é datado a aproximadamente 2800 a.C \cite{fogaca_historia_sabao}. Onde, os sumérios, povo que vivia na região da Babilônia, usava misturas de gordura animal e cinzas vegetais para aplicações como limpeza e tratamento de infecções \cite{rittner1995sabao}. Essa mesma técnica para produção desses materiais foi mencionada em diversas escritas romanas, sua adoção pelo povo romano incidi num maior uso e possibilita a melhor elaboração do mesmo \cite{barbosa1995xampus}.

A origem da palavra "sabão" ~remonta ao termo latino \textit{sapo}, relacionado a uma antiga narrativa da qual sua produção teria surgido nas encostas do Monte Sapo. Nesse local, resíduos gordurosos provenientes de sacrifícios de animais escorriam pela face do monte e se combinavam com cinzas geradas por fogueiras, dando origem a uma substância com propriedades de limpeza \cite{shreve1956chemical}. Note que tal local  "Monte Sapo" é estreitamente lendário, muito provavelmente foi criado simplesmente para explicar a origem
do sabão. 

Durante a Idade Média, a fabricação do sabão consolidou-se como uma importante atividade artesanal em relevantes centros urbanos europeus, destacando-se os polos produtores de Veneza e Marselha, reconhecidos pelo emprego de azeite de oliva e de compostos alcalinos obtidos a partir de vegetais \cite{rittner1995sabao}.

Um marco muito importante para a história do sabão foi a sintetização de carbonato de sódio (barrilha) a partir do sal de cozinha, iniciativa tomada Nicolas Leblanc no século XVIII. Dessa forma, possibilitando a produção desse material em média escala. Logo em seguida, Michel Eugene Chevreul, melhorou ainda mais a formulação, estabelecendo os mecanismos básicos da saponificação \cite{selinger1986chemistry}.

A Revolução Industrial, ao longo do século XIX, transformou profundamente o status do sabão, antes restrito  às classes mais ricas da sociedade, o produto tornou-se acessível à população em geral, ao passo que seu processo produtivo passou a incorporar fragrâncias e aditivos diversos \cite{rittner1995sabao}. Na contemporaneidade, a produção de sabão se distribui entre o âmbito artesanal e o industrial, respondendo a distintas demandas de consumo e integrando tanto avanços tecnológicos quanto considerações de ordem ambiental em sua cadeia de fabricação \cite{castro2020producao}.